\abstract{Alternating left-right auditory stimulation is sometimes discussed alongside bilateral stimulation used in eye movement desensitization and reprocessing (EMDR), but this does not establish a specific therapeutic mechanism for attention-deficit/hyperactivity disorder (ADHD). This companion paper therefore removes bilateral audio from the core ARC-ADHD intervention and treats it as a separate proof-of-mechanism experiment. The primary contrast is alternating left-right auditory stimulation versus acoustically matched non-alternating stimulation, with a quiet regulation condition included to estimate nonspecific auditory effects. The study targets proximal recovery after a controlled challenge rather than ADHD symptom reduction. Candidate mechanisms include attentional orienting, temporal pacing, dual-task interference, and nonspecific rhythmic regulation. These mechanisms make different predictions and should not be collapsed into an EMDR analogy. The design uses within-participant randomized eligible events, prespecified acoustic parameters, tolerability limits, washout rules, and a falsification criterion requiring alternating stimulation to outperform the matched auditory control. A null result against the matched control would argue against retaining a specifically bilateral mechanism in ARC-ADHD.}
